\documentclass[12pt]{article}
\usepackage{latexblog}

% ============================================================
% Blog Metadata
% ============================================================
\blogtitle{LaTeX(3)：使用TikZ绘制图形}
\blogdate{2020-02-09}
\blogtags{TikZ, 其他, TeX, ~LaTeX入门, LaTeX}
\bloglang{zh}
\blogsummary{}

\begin{document}
\makeblogtitle

\href{https://www.ctan.org/pkg/pgf}{PGF}是一个用来进行图形绘制的（底层）包，TikZ是利用这个包实现的用户友好的接口。所以通常在LaTeX中会用TikZ来进行矢量图形的绘制。

\section{基本概念}\label{ux57faux672cux6982ux5ff5}

\subsection{基本语法}\label{ux57faux672cux8bedux6cd5}

要使用tikz进行图形绘制只需要简单引入tikz宏包，并将绘制代码包含在一个上下文中就可以了:

\begin{Shaded}
\begin{Highlighting}[]
\BuiltInTok{\textbackslash{}usepackage}\NormalTok{\{}\ExtensionTok{tikz}\NormalTok{\}}

\KeywordTok{\textbackslash{}begin}\NormalTok{\{}\ExtensionTok{tikzpicture}\NormalTok{\}}
\CommentTok{\% ...}
\KeywordTok{\textbackslash{}end}\NormalTok{\{}\ExtensionTok{tikzpicture}\NormalTok{\}}
\end{Highlighting}
\end{Shaded}

同样如果希望把图片作为一个figure，那么再套一层:

\begin{Shaded}
\begin{Highlighting}[]
\BuiltInTok{\textbackslash{}usepackage}\NormalTok{\{}\ExtensionTok{graphicx}\NormalTok{\}}
\BuiltInTok{\textbackslash{}usepackage}\NormalTok{\{}\ExtensionTok{tikz}\NormalTok{\}}

\KeywordTok{\textbackslash{}begin}\NormalTok{\{}\ExtensionTok{document}\NormalTok{\}}
    \KeywordTok{\textbackslash{}begin}\NormalTok{\{}\ExtensionTok{figure}\NormalTok{\}[h]}
    \KeywordTok{\textbackslash{}begin}\NormalTok{\{}\ExtensionTok{tikzpicture}\NormalTok{\}}
        \FunctionTok{\textbackslash{}draw}\NormalTok{ (0, 0) {-}{-} (1, 1);}
    \KeywordTok{\textbackslash{}end}\NormalTok{\{}\ExtensionTok{tikzpicture}\NormalTok{\}}
    \FunctionTok{\textbackslash{}caption}\NormalTok{\{这是一条直线\}}
    \KeywordTok{\textbackslash{}end}\NormalTok{\{}\ExtensionTok{figure}\NormalTok{\}}
\KeywordTok{\textbackslash{}end}\NormalTok{\{}\ExtensionTok{document}\NormalTok{\}}
\end{Highlighting}
\end{Shaded}

\subsection{坐标系和缩放}\label{ux5750ux6807ux7cfbux548cux7f29ux653e}

默认情况下TikZ使用的是笛卡尔坐标系，即是这样：

另外，TikZ提供了一个缩放的选项，可以用来缩放图形，所以不必要担心绝对坐标的问题：

\begin{Shaded}
\begin{Highlighting}[]
\KeywordTok{\textbackslash{}begin}\NormalTok{\{}\ExtensionTok{tikzpicture}\NormalTok{\}[scale=0.5]}
\end{Highlighting}
\end{Shaded}

甚至还可以针对x轴和y轴分别设置缩放：

\begin{Shaded}
\begin{Highlighting}[]
\KeywordTok{\textbackslash{}begin}\NormalTok{\{}\ExtensionTok{tikzpicture}\NormalTok{\}[xscale=0.5, yscale=0.3]}
\end{Highlighting}
\end{Shaded}

在TikZ中，默认的单位是厘米(cm)。如果希望改变这个值，可以这样设置：

\begin{Shaded}
\begin{Highlighting}[]
\KeywordTok{\textbackslash{}begin}\NormalTok{\{}\ExtensionTok{tikzpicture}\NormalTok{\}[x=2cm,y=1.5cm]}

\CommentTok{\% 或者}
\KeywordTok{\textbackslash{}begin}\NormalTok{\{}\ExtensionTok{tikzpicture}\NormalTok{\}[x=\{(2cm,0cm)\},y=\{(0cm,1.5cm)\}]}
\end{Highlighting}
\end{Shaded}

\section{图形绘制}\label{ux56feux5f62ux7ed8ux5236}

\subsection{直线}\label{ux76f4ux7ebf}

绘制直线：

\begin{Shaded}
\begin{Highlighting}[]
\FunctionTok{\textbackslash{}draw}\NormalTok{ (0, 0) {-}{-} (1, 1);}
\end{Highlighting}
\end{Shaded}

绘制折线：

\begin{Shaded}
\begin{Highlighting}[]
\FunctionTok{\textbackslash{}draw}\NormalTok{ (0, 0) {-}{-} (1, 1) {-}{-} (2, 2) {-}{-} (1, 0) {-}{-} (0, 3);}
\end{Highlighting}
\end{Shaded}

绘制背景网格:

\begin{Shaded}
\begin{Highlighting}[]
\FunctionTok{\textbackslash{}draw}\NormalTok{[help lines] (0,0) grid (3,3);}
\end{Highlighting}
\end{Shaded}

\subsubsection{箭头}\label{ux7badux5934}

如果希望绘制箭头也十分方便：

\begin{Shaded}
\begin{Highlighting}[]
\FunctionTok{\textbackslash{}draw}\NormalTok{ [{-}\textgreater{}] (0,0) {-}{-} (2,0);        }\CommentTok{\%→}
\FunctionTok{\textbackslash{}draw}\NormalTok{ [\textless{}{-}] (0, {-}0.5) {-}{-} (2,{-}0.5); }\CommentTok{\%←}
\FunctionTok{\textbackslash{}draw}\NormalTok{ [|{-}\textgreater{}] (0,{-}1) {-}{-} (2,{-}1);     }\CommentTok{\%带尾巴的箭头}
\FunctionTok{\textbackslash{}draw}\NormalTok{ [\textless{}{-}\textgreater{}] (0, 0) {-}{-} (1, 1);     }\CommentTok{\%双向箭头}
\end{Highlighting}
\end{Shaded}

\subsubsection{线的粗细}\label{ux7ebfux7684ux7c97ux7ec6}

线的粗细可以用如下来表示：

\begin{Shaded}
\begin{Highlighting}[]
\FunctionTok{\textbackslash{}draw}\NormalTok{ [ultra thin] (0, 1) {-}{-} (2, 1)}
\end{Highlighting}
\end{Shaded}

总共可用的粗细如下：

或者直接指定线的粗细，默认的单位是点:

\begin{Shaded}
\begin{Highlighting}[]
\FunctionTok{\textbackslash{}draw}\NormalTok{ [line width=12] (0,0) {-}{-} (2,0);}
\FunctionTok{\textbackslash{}draw}\NormalTok{ [line width=0.2cm] (4,.75) {-}{-} (5,.25);}
\end{Highlighting}
\end{Shaded}

除此之外另一个选项是\texttt{{[}help\ lines{]}}，用来绘制灰色的参考线。

\begin{Shaded}
\begin{Highlighting}[]
\FunctionTok{\textbackslash{}draw}\NormalTok{ [help lines] (0, 5) {-}{-} (0, 0) {-}{-} (5, 0);}
\FunctionTok{\textbackslash{}draw}\NormalTok{ [line width=2pt] (0, 0) {-}{-} (5, 5);}
\FunctionTok{\textbackslash{}draw}\NormalTok{ [very thin] (0, 3) {-}{-} (4, 0);}
\FunctionTok{\textbackslash{}draw}\NormalTok{ [thin] (0, 2) {-}{-} (5, 2);}
\end{Highlighting}
\end{Shaded}

\subsubsection{样式及颜色}\label{ux6837ux5f0fux53caux989cux8272}

样式可以分为：

\begin{itemize}
\tightlist
\item
  虚线 \texttt{\textbackslash{}draw\ {[}dashed{]}}
\item
  实线 \texttt{\textbackslash{}draw\ {[}dotted{]}}
\end{itemize}

颜色有很多直接可以用的颜色表示，类似css一样：

\begin{quote}
red, green, blue, cyan, magenta, yellow, black, gray, darkgray, lightgray,brown, lime, olive, orange, pink, purple, teal, violetand white
\end{quote}

一个较为完整的例子:

\begin{Shaded}
\begin{Highlighting}[]
\KeywordTok{\textbackslash{}begin}\NormalTok{\{}\ExtensionTok{tikzpicture}\NormalTok{\}}
    \FunctionTok{\textbackslash{}draw}\NormalTok{ [help lines] (0, 5) {-}{-} (0, 0) {-}{-} (5, 0) node [right=3]\{Nice sample!\};}
    \FunctionTok{\textbackslash{}draw}\NormalTok{ [dashed, red, line width=2pt] (0, 0) {-}{-} (5, 5);}
    \FunctionTok{\textbackslash{}draw}\NormalTok{ [blue, very thin] (0, 3) {-}{-} (4, 0);}
    \FunctionTok{\textbackslash{}draw}\NormalTok{ [dotted, thin] (0, 2) {-}{-} (5, 2);}
\KeywordTok{\textbackslash{}end}\NormalTok{\{}\ExtensionTok{tikzpicture}\NormalTok{\}}
\end{Highlighting}
\end{Shaded}

\subsection{几何图形}\label{ux51e0ux4f55ux56feux5f62}

\subsubsection{矩形}\label{ux77e9ux5f62}

\begin{Shaded}
\begin{Highlighting}[]
\FunctionTok{\textbackslash{}draw}\NormalTok{ [blue] (0,0) rectangle (1.5, 1);}
\end{Highlighting}
\end{Shaded}

\subsubsection{网格}\label{ux7f51ux683c}

\begin{Shaded}
\begin{Highlighting}[]
\FunctionTok{\textbackslash{}draw}\NormalTok{ [blue] (0,0) grid (1.5, 1);}
\end{Highlighting}
\end{Shaded}

\subsubsection{圆}\label{ux5706}

\begin{Shaded}
\begin{Highlighting}[]
\FunctionTok{\textbackslash{}draw}\NormalTok{ [red, ultra thick] (3, 0.5) circle [radius=0.5];}
\end{Highlighting}
\end{Shaded}

\subsubsection{弧线}\label{ux5f27ux7ebf}

\begin{Shaded}
\begin{Highlighting}[]
\FunctionTok{\textbackslash{}draw}\NormalTok{ [gray, ultra thick] (6,0) arc [radius=1, start angle=45, end angle= 120];}
\end{Highlighting}
\end{Shaded}

这个弧线的表示方法比较有意思，代表从(6, 0)出发，半径为1，初始角度为45°，当变成120°的时候停止。另一种方法：

\begin{Shaded}
\begin{Highlighting}[]
\FunctionTok{\textbackslash{}draw}\NormalTok{[very thick] (0,0) to [out=90,in=195] (2,1.5);}
\end{Highlighting}
\end{Shaded}

表示从(0,0)开始， 到(2, 1.5)这个点，起始角度为90°，到达的角度为195°。

感觉很难控制这个\ldots.

\subsubsection{圆角折线}\label{ux5706ux89d2ux6298ux7ebf}

加多一个\texttt{rounded\ corners}就可以把折线变成圆角的了：

\begin{Shaded}
\begin{Highlighting}[]
\FunctionTok{\textbackslash{}draw}\NormalTok{ [\textless{}{-}\textgreater{}, rounded corners, thick, purple] (0,2) {-}{-} (0,0) {-}{-} (3,0);}
\end{Highlighting}
\end{Shaded}

\begin{itemize}
\tightlist
\item
  \href{https://www.latex-tutorial.com/tutorials/tikz/}{Draw pictures in LaTeX - With tikz/pgf}
\item
  \href{http://cremeronline.com/LaTeX/minimaltikz.pdf}{A very minimal introduction to TikZ∗}
\end{itemize}


\end{document}
